
Sistemas neuro-fuzzy evolutivos representam o estado da arte para o tratamento de problemas complexos e dinâmicos, como a detecção de fraudes em grandes bases de dados. Ao combinar a capacidade adaptativa dos sistemas evolutivos, a modelagem de incertezas dos sistemas fuzzy e o poder de aprendizagem das redes neurais, esses modelos conseguem ajustar suas regras, funções de pertinência e parâmetros de decisão de forma contínua, acompanhando a evolução dos padrões de fraude. Essa abordagem permite alta sensibilidade e precisão mesmo diante de estratégias fraudulentas inéditas, tornando-se uma solução robusta, flexível e interpretável para ambientes adversariais e em constante transformação.

\section{Etapas da Metodologia}

\begin{enumerate}
\item \textbf{Revisão Bibliográfica e Fundamentação Teórica}
\begin{itemize}
\item Realizar levantamento sistemático da literatura sobre sistemas inteligentes evolutivos, incluindo trabalhos pioneiros e desenvolvimentos recentes
\item Analisar metodologias de detecção de anomalias e fraudes em sistemas complexos
\item Estudar algoritmos de aprendizado incremental e online, com foco em adaptação estrutural e paramétrica
\item Investigar técnicas de tratamento de drift conceitual e adaptação a ambientes não-estacionários
\item Identificar lacunas teóricas e práticas que possam ser exploradas na proposição de novos algoritmos
\end{itemize}

\item \textbf{Proposição e Implementação de Algoritmos}
\begin{itemize}
\item Propor novos algoritmos baseados em sistemas fuzzy evolutivos adaptados ao contexto de detecção de fraudes
\item Desenvolver mecanismos de adaptação estrutural automática (criação, fusão e remoção de regras/grupos)
\item Implementar técnicas de aprendizado incremental para processamento de dados em tempo real
\item Integrar métodos de pertinência fuzzy para tratamento de incertezas e modelagem de padrões ambíguos
\item Garantir interpretabilidade através de representação linguística e explicabilidade das decisões
\item Implementar algoritmos em linguagem Python, priorizando código modular, reutilizável e documentado
\item Desenvolver dashboards interativos para visualização dos resultados, utilizando bibliotecas Python como Dash ou Streamlit, de modo a facilitar a análise exploratória, a interpretação dos dados e a comunicação dos achados
\item Implementar um sistema web para disponibilização dos resultados, relatórios e ferramentas de análise, permitindo acesso remoto, interação com os dados e compartilhamento dos achados com diferentes públicos de interesse
\end{itemize}

\item \textbf{Validação Experimental e Análise Comparativa}
\begin{itemize}
\item Aplicar algoritmos propostos em datasets reais de sistemas agropecuários
\item Comparar desempenho com métodos estado da arte usando métricas padronizadas
\item Realizar análises estatísticas de significância dos resultados
\item Avaliar capacidade de adaptação em cenários de mudança conceitual
\item Analisar eficiência computacional e escalabilidade dos métodos propostos
\end{itemize}
\end{enumerate}


\section{Cronograma de Execução}

O planejamento temporal da pesquisa foi elaborado de modo a garantir a execução ordenada e eficiente de todas as etapas previstas, desde a fundamentação teórica até a validação dos resultados. O cronograma detalhado, incluindo as atividades de pesquisa, desenvolvimento, análise e as disciplinas obrigatórias do doutorado, encontra-se apresentado no Capítulo~\ref{chap:cronograma}. Ressalta-se que o acompanhamento das disciplinas cursadas, a cursar e aproveitáveis está centralizado naquele capítulo, evitando redundâncias neste.


\section{Aspectos Éticos}

Todos os procedimentos da pesquisa respeitarão rigorosamente os princípios éticos aplicáveis ao uso de dados sensíveis, em conformidade com a legislação vigente. Serão adotadas as seguintes diretrizes:
\begin{itemize}
    \item Anonimização e proteção integral dos dados utilizados;
    \item Observância à Lei Geral de Proteção de Dados (LGPD);
    \item Obtenção de autorizações institucionais formais para uso dos dados;
    \item Garantia de que nenhuma informação permita a identificação de indivíduos, propriedades ou organizações.
\end{itemize}


\section{Métricas de Avaliação}

Para garantir a robustez e a comparabilidade dos resultados, a avaliação dos algoritmos propostos será realizada por meio de métricas consagradas na literatura de detecção de fraudes e sistemas evolutivos. Serão consideradas métricas de precisão, classificação e testes estatísticos:


\subsection{Métricas de Precisão}
As métricas de precisão são utilizadas para avaliar o desempenho dos modelos em tarefas de previsão de valores contínuos:
\begin{itemize}
    \item \textbf{Raiz do Erro Quadrático Médio (RMSE):} Mede o erro médio das previsões realizadas pelo modelo em relação aos valores reais. Quanto menor o RMSE, melhor o desempenho do modelo, pois indica que as previsões estão mais próximas dos valores observados.
    \[RMSE = \sqrt{\frac{\sum_{t=1}^{N}(y_t - \hat{y}_t)^2}{N}}\]
    Onde $y_t$ é o valor real, $\hat{y}_t$ é o valor previsto e $N$ é o número total de observações.
    \item \textbf{Índice de Erro Não Dimensional (NDEI):} É a razão entre o RMSE e o desvio padrão dos valores previstos. Permite comparar o erro relativo ao grau de dispersão dos dados, tornando a métrica adimensional e facilitando a comparação entre diferentes conjuntos de dados.
    \[NDEI = \frac{RMSE}{\sigma(\hat{y}_t)}\]
    Onde $\sigma(\hat{y}_t)$ é o desvio padrão das previsões.
    \item \textbf{Soma Residual dos Quadrados (RSS):} Representa a soma dos quadrados das diferenças entre os valores reais e previstos. É utilizada para avaliar o ajuste do modelo: quanto menor o RSS, melhor o ajuste.
    \[RSS = \sum_{t=1}^{N}(y_t - \hat{y}_t)^2\]
\end{itemize}


\subsection{Métricas de Classificação}
As métricas de classificação são empregadas para avaliar o desempenho dos modelos em tarefas de detecção de fraudes (classificação binária):
\begin{itemize}
    \item \textbf{Precisão (Precision):} Indica a proporção de casos classificados como positivos que realmente são positivos. É útil para avaliar o desempenho do modelo em situações onde o custo de falsos positivos é alto.
    \[\text{Precisão} = \frac{TP}{TP + FP}\]
    Onde $TP$ são os verdadeiros positivos e $FP$ os falsos positivos.
    \item \textbf{Recall (Sensibilidade):} Mede a capacidade do modelo de identificar corretamente todos os casos positivos. É importante em contextos onde a detecção de todos os casos reais é prioritária.
    \[\text{Recall} = \frac{TP}{TP + FN}\]
    Onde $FN$ são os falsos negativos.
    \item \textbf{F1-score:} É a média harmônica entre precisão e recall, fornecendo uma medida única que equilibra ambos os aspectos. É especialmente útil quando há desequilíbrio entre as classes.
    \[F1 = 2 \cdot \frac{\text{Precisão} \cdot \text{Recall}}{\text{Precisão} + \text{Recall}}\]
    \item \textbf{Acurácia:} Representa a proporção de predições corretas (tanto positivas quanto negativas) em relação ao total de casos avaliados.
    \[\text{Acurácia} = \frac{TP + TN}{TP + TN + FP + FN}\]
    Onde $TN$ são os verdadeiros negativos.
    \item \textbf{Área sob a Curva ROC (AUC-ROC):} Mede a capacidade do classificador em distinguir entre as classes. Quanto mais próximo de 1, melhor o desempenho do modelo. Um valor de 0,5 indica desempenho aleatório.
\end{itemize}

\subsection{Testes Estatísticos}
Os testes estatísticos são fundamentais para validar se as diferenças observadas entre os modelos são estatisticamente significativas e não fruto do acaso. No contexto deste projeto, será utilizado o seguinte teste:
\begin{itemize}
    \item \textbf{Teste Morgan-Granger-Newbold (MGN):} Utilizado para comparar o desempenho de dois modelos preditivos, avaliando se a diferença entre os erros (resíduos) de ambos é estatisticamente significativa. O teste calcula um valor baseado na correlação entre os resíduos dos modelos, permitindo verificar se um modelo apresenta desempenho superior ao outro de forma consistente.
    \[MGN = \hat{\rho}_e \sqrt{\frac{N-1}{1-\hat{\rho}_e^2}}\]
    Onde $\hat{\rho}_e$ é a correlação entre os resíduos dos modelos comparados e $N$ é o número de observações. Valores elevados de MGN indicam maior diferença estatística entre os modelos.
\end{itemize}

\section{Base de Dados}

O estudo utilizará dados reais do sistema de gerenciamento agropecuário, especificamente registros de Guias de Trânsito Animal (GTA) coletados ao longo dos anos. A base de dados principal possui as seguintes características:

\subsection{Estrutura dos Dados}
O dataset contém as seguintes variáveis principais:

\begin{table}[h]
\centering
\caption{Estrutura da Base de Dados de GTAs}
\begin{tabular}{|l|p{8cm}|}
\hline
\textbf{Variável} & \textbf{Descrição} \\
\hline
nr\_gta & Número identificador único da Guia de Trânsito Animal \\
\hline
nr\_serie & Número de série da GTA \\
\hline
id\_gta & Identificador interno da GTA no sistema \\
\hline
dt\_emissao\_gta & Data de emissão da guia \\
\hline
cod\_produtor\_origem & Código do produtor de origem dos animais \\
\hline
cod\_produtor\_destino & Código do produtor de destino dos animais \\
\hline
id\_tipo\_finalidade\_gta & Identificador do tipo de finalidade da movimentação \\
\hline
ds\_meio\_transporte & Descrição do meio de transporte utilizado \\
\hline
cod\_municipio\_origem & Código do município de origem \\
\hline
nome\_municipio\_origem & Nome do município de origem \\
\hline
estado\_origem & Estado de origem (MG) \\
\hline
cod\_municipio\_destino & Código do município de destino \\
\hline
nome\_municipio\_destino & Nome do município de destino \\
\hline
estado\_destino & Estado de destino \\
\hline
qtd & Quantidade de animais transportados \\
\hline
\end{tabular}
\end{table}

\subsection{Potencial para Detecção de Fraudes}
Esta base de dados apresenta características ideais para o desenvolvimento e validação de algoritmos de detecção de fraudes:

\begin{itemize}
\item \textbf{Volume significativo:} Com mais de 8 milhões de registros, permite análises estatísticas robustas
\item \textbf{Diversidade temporal:} Dados históricos possibilitam análise de padrões evolutivos
\item \textbf{Informações relacionais:} Conexões entre produtores, municípios e finalidades permitem análise de redes
\item \textbf{Variáveis de controle:} Dados de quantidade, datas e localizações facilitam identificação de anomalias
\end{itemize}


\subsection{Tratamento e Preparação dos Dados}
Os dados serão submetidos aos seguintes processos:
\begin{enumerate}
    \item \textbf{Limpeza e validação:} Identificação e tratamento de valores ausentes, duplicatas e inconsistências
    \item \textbf{Anonimização:} Remoção ou codificação de informações que possam identificar produtores específicos
    \item \textbf{Engenharia de características:} Criação de variáveis derivadas para análise de padrões (frequência de movimentação, distâncias, sazonalidade)
    \item \textbf{Normalização:} Padronização de escalas para algoritmos de aprendizado de máquina
    \item \textbf{Detecção de outliers:} Identificação preliminar de transações atípicas
\end{enumerate}

\subsection{Proposta de Sistema Inteligente Fuzzy Evolutivo}

Com base na base de dados de GTAs, propõe-se o desenvolvimento de um sistema inteligente fuzzy evolutivo para detecção de fraudes em movimentações agropecuárias. O sistema seguirá as seguintes etapas:
\begin{enumerate}
    \item \textbf{Pré-processamento dos Dados:} Limpeza, anonimização e normalização das variáveis, além da criação de atributos derivados, como frequência de movimentação, distâncias e sazonalidade.
    \item \textbf{Modelagem Fuzzy Evolutiva:} Definição de variáveis linguísticas (ex: quantidade baixa/média/alta, distância curta/longa, frequência incomum), construção de funções de pertinência para representar incertezas e inicialização de regras fuzzy para padrões normais e suspeitos.
    \item \textbf{Evolução Online:} O sistema ajusta automaticamente as funções de pertinência e as regras fuzzy conforme novos dados são inseridos, utilizando algoritmos evolutivos para criação, fusão e remoção de regras. O aprendizado incremental permite adaptação contínua a novos comportamentos e fraudes.
    \item \textbf{Detecção de Fraudes:} Para cada nova GTA registrada, o sistema calcula o grau de pertinência a padrões normais e suspeitos, sinalizando casos anômalos para análise detalhada.
    \item \textbf{Visualização e Interação:} Dashboards interativos apresentam estatísticas, alertas e padrões detectados, facilitando a tomada de decisão por gestores e auditores.
\end{enumerate}

Essa abordagem permite identificar fraudes mesmo em cenários dinâmicos e incertos, aproveitando o potencial adaptativo dos sistemas fuzzy evolutivos e a riqueza de informações do banco de dados agropecuário.

Portanto, a metodologia proposta permite alcançar os objetivos gerais e específicos do projeto. Cada etapa está alinhada com os objetivos previamente definidos, desde a análise do problema e coleta de dados até a validação do modelo e proposição de melhorias. Isso garante que a pesquisa seja conduzida de forma sistemática, transparente e orientada a resultados teóricos, práticos e aplicáveis.
